\centerline{\bf \Huge Тренировочная олимпиада II}
\centerline{\bf \Huge }

\begin{flushright}
        \scriptsize{}
\end{flushright}

\problem{1}
Дима и Саша весь день катались на метро и автобусе. Дима ездил по карте «Тройка», а Саша по карте «Единый». Цена проезда по карте «Тройка» на автобусе составляет 37 рублей, а на метро - 48 рублей. Цена проезда по карте «Единый» на автобусе и на метро одинакова и составляет 41 рубль.

Известно, что за день они совершили одинаковое количество поездок и потратили одинаковую сумму, меньшую 500 рублей. Сколько потратил один из них?


\problem{2}
Ювелир изготовил 6 одинаковых по виду серебряных украшений массой 22 г, 23 г, 24 г, 32 г, 34 г и 36 г и поручил своему подмастерью выбить на каждом украшении его массу. Может ли ювелир за два взвешивания на чашечных весах без стрелок и гирек определить, не перепутал ли подмастерье украшения?

\problem{3}
На сколько равных восьмиугольников можно разрезать квадрат размером $8 \times 8$ ? (Все разрезы должны проходить по линиям сетки.)

\problem{4}
На сторонах $A B$ и $B C$ треугольника $A B C$ отметили точки $P$ и $Q$ соответственно. Известно, что

$$
B P=B C, \quad P Q=A C=21, \quad Q C=12, \quad \angle B A C=\angle P Q C=60^{\circ} .
$$


Найдите длину отрезка $A P$.

\problem{5}
Все натуральные делители числа $n$ выписали в ряд: $d_1, d_2, \ldots, d_k$. Известно, что $d_1+$ $d_2+\ldots+d_k=4368$ и $\frac{1}{d_1}+\frac{1}{d_2}+\ldots+\frac{1}{d_k}=\frac{52}{15}$. Найдите, чему равно $n$.